%% kapa: exam=1 round=25 year=2014 period=1 form=A subject=economics subject_ko=경제원론 seq=8 number=48
%% source: Q-Net 감정평가사 기출문제 (cst00309) artlSeq=5101207 file=extracted/1008/5101207/2014년 제25회 감정평가사 1차시험 기출문제/1교시(A형).hwp
%% conversion: kordoc 4.24.51 -> tools/gen1.py (auto); review=proofread-sonnet; flags=none
\begin{question}{48}
이윤을 극대화하는 기업 A와 B의 생산량과 이윤행렬은 다음과 같다. A는 슈타켈버그(Stackelberg) 모형의 선도자, B는 추종자로 행동할 때 A와 B의 생산량($Q_{ A }$, $Q_{ B }$)은? (단, 이윤행렬의 괄호안의 수에서 왼쪽은 A의 이윤이고, 오른쪽은 B의 이윤이다.) \par
{\small\centering
\begin{tabular}{|c|c|c|c|c|}
\hline
\multicolumn{2}{|c|}{\multirow{2}{*}{}} & \multicolumn{3}{|c|}{\makecell{기업 B의 생산량}} \\
\cline{3-5}
\multicolumn{2}{|c|}{} & \makecell{15} & \makecell{20} & \makecell{30} \\
\hline
\multirow{3}{*}{\makecell{기업 A의 생산량}} & \makecell{15} & \makecell{(450, 450)} & \makecell{(375, 500)} & \makecell{(225, 450)} \\
\cline{2-5}
 & \makecell{20} & \makecell{(500, 375)} & \makecell{(400, 400)} & \makecell{(200, 300)} \\
\cline{2-5}
 & \makecell{30} & \makecell{(450, 225)} & \makecell{(300, 200)} & \makecell{(0, 0)} \\
\hline
\end{tabular}\par}

\begin{choices}
  \item (15, 15)
  \item (20, 15)
  \item (20, 20)
  \item (30, 15)
  \item (30, 20)
\end{choices}
\end{question}
\answer{4}
